% Úvod a metodika případových studií
% Poznámka: níže uvedená metodika obsahuje praktická doporučení; část obsahu je obecné metodické doporučení (general background knowledge).

Tato sekce vysvětluje, jak jsou jednotlivé případové studie strukturovány a jak je číst tak, aby byly použitelné při redesignu předmětů na VUT. Hlavním cílem je poskytnout čtenáři jasný checklist pro posouzení přenositelnosti jednotlivého příkladu do lokálního kontextu a krok‑za‑krokem pracovní postup pro pilotní adaptaci (general background knowledge).

\begin{itemize}
  \item Struktura každé případové studie (pole, která vždy najdete):
    \begin{itemize}
      \item Instituce a kontext kurzu (velikost kurzu, formát: prezenčně/online/hybrid).
      \item Cíl a pedagogické metriky (co se měřilo).
      \item AI řešení (co bylo použito: nástroj, konfigurace, integrace).
      \item Implementace (fáze nasazení, role lidí, technická integrace).
      \item Výsledky a evidence (kvantitativní a kvalitativní poznatky).
      \item Lessons learned (co fungovalo/na co si dát pozor).
      \item Odkazy a zdroje pro další čtení.
    \end{itemize}
  \item Jak číst kazuistiku — rychlý průvodce pro garanta předmětu:
    \begin{enumerate}
      \item Identifikujte, které části řešení jsou technické (model, licence, integrace) a které pedagogické (úkoly, hodnocení, procesní důkazy).
      \item Hledejte konkrétní měřitelné výsledky nebo indikátory úspěchu, které lze reprodukovat ve vašem předmětu (např. čas na opravu, změna engagementu apod.) (general background knowledge).
      \item Zmapujte rizika a mitigace popsané v kazuistice (data, soukromí, akademická integrita).
      \item Zvažte nutné kroky pro lokalizaci: překlad, adaptace zadání na lokální data, ověření licencí a kompatibility s IT prostředím VUT.
    \end{enumerate}
\end{itemize}

Praktické ukázky nástrojů v kazuistikách často obsahují příklady edukačních aplikací a vzdělávacích akcelerátorů; při posuzování jejich přenositelnosti věnujte pozornost, zda řešení vyžaduje školní licence nebo specifické integrační kroky. Jako ilustraci typu nástrojů, které se v praxi používají, lze uvést Minecraft Education pro didaktické světy a Hodinu kódu (viz popis využití v materiálu) \cite{kb_ai_101_tipu_pdf_04e4a115_page_15_2}. Podobně jsou v pedagogických scénářích zmíněny vzdělávací akcelerátory pro procvičování čtení a matematiky (funkce „Pokrok ve čtení“ a „Pokrok v matematice“) — tyto příklady ukazují, jak může systém předzpracovat a vyhodnotit studentovy odpovědi a navrhnout personalizaci \cite{kb_ai_101_tipu_pdf_04e4a115_page_7_1, kb_ai_101_tipu_pdf_04e4a115_page_8_1}.
  
Krátký checklist pro posouzení přenositelnosti případu:
\begin{itemize}
  \item Je cíl pedagogicky kompatibilní s očekávanými výstupy předmětu? (general background knowledge)
  \item Jaké jsou technické požadavky a jsou splnitelné v prostředí VUT (LMS, SSO, enterprise licence)? (general background knowledge)
  \item Jaké údaje se zpracovávají a je potřeba provést DPIA / DPA před nasazením? (general background knowledge)
  \item Existuje jasný plán pilotu (malý rozsah → validace → škálování) a metrika úspěchu? (general background knowledge)
  \item Jsou ve studii popsané konkrétní „lessons learned“, které eliminují známé rizikové kroky? (general background knowledge)
\end{itemize}

Jak využít kazuistiku při redesignu předmětu — doporučený postup (3 kroky, praktické):
\begin{enumerate}
  \item Extrahujte z kazuistiky komponenty, které chcete replikovat: (a) typ úkolu, (b) role AI, (c) požadavek na důkazy studentovy práce.
  \item Navrhněte malý pilot v rámci jednoho tématu/sem. hodiny: definujte vstupy, výstupy, metriky a lidské validační body (např. pedagog prověří 10\% výstupů).
  \item Vyhodnoťte pilot podle předem definovaných metrik a zdokumentujte „lessons learned“ pro širší nasazení nebo úpravu sylabu (general background knowledge).
\end{enumerate}

Na závěr: při čtení kapitol s případovými studiemi hledejte konkrétní technické a pedagogické body, které lze vyzkoušet v malém měřítku; využijte přiložené checklisty této kapitoly jako šablonu pro vlastní pilotní plán (general background knowledge).